
\begin{Info}
% Universidade
{UNIVERSIDADE DO VALE DO ITAJAÍ}
% Escola
{ESCOLA POLITÉCNICA}
% Curso
{CURSO DE CIÊNCIA DA COMPUTAÇÃO}
% Titulo
% TODO: validar com o orientador. Alternativas em discussao:
%   (A) ...COM EXTRACAO VIA LLM
%   (B) ...COM EXTRACAO ESTRUTURADA VIA LLM
%   (C) ...COM PIPELINE SERVERLESS E EXTRACAO VIA LLM
{ROLÊ SC: APLICATIVO MOBILE AGREGADOR DE EVENTOS CULTURAIS E MUNICIPAIS DO LITORAL NORTE DE SANTA CATARINA COM EXTRAÇÃO AUTOMATIZADA VIA LLM}
% Autor
{Gabriel Sodré da Silva Prates}
% Cidade e Data
{Itajaí (SC), novembro de 2026}
% Nome da Área de concentração
{Engenharia de Software}
% Orientador(a)
{Alex Luciano Roesler Rese, Me.}
{}
\end{Info}


\begin{Resumo}
A descoberta de eventos culturais, esportivos, gastronômicos e comunitários no litoral norte de Santa Catarina enfrenta um problema de fragmentação: cada prefeitura, fundação cultural e produtor publica suas programações em sites próprios e redes sociais, sem padronização nem agregação regional. Este trabalho de conclusão de curso (TCC) apresenta o desenvolvimento do \textit{Rolê SC}, um aplicativo móvel agregador de eventos dos municípios de Itajaí, Camboriú, Balneário Camboriú, Porto Belo, Itapema e Bombinhas, sustentado por um \textit{pipeline serverless} que automatiza a coleta e a estruturação dessas informações. O ecossistema é composto por quatro componentes independentes implantados na nuvem da \textit{Amazon Web Services}: um \textit{web crawler} baseado em \textit{Playwright} que coleta páginas dos sites oficiais; um processador de linguagem natural que utiliza o modelo \textit{Llama 3.1 8B Instruct}, servido pela plataforma \textit{Groq}, para extrair eventos estruturados ou descartar páginas que não correspondem a eventos; uma \textit{API REST} em \textit{TypeScript} sobre o \textit{framework Fastify}, autenticada por \textit{JWT}; e um aplicativo móvel Android construído em \textit{React Native} com \textit{Expo}. O trabalho contempla também o estudo do modelo de negócio, a análise de concorrentes, o levantamento de requisitos e a documentação do sistema. Nesta primeira etapa (TCC-II), foram implementados o \textit{pipeline} de coleta e extração para os cinco municípios-alvo, os fluxos de cadastro, verificação de email, autenticação, recuperação e alteração de senha, edição de perfil e \textit{upload} de avatar, bem como a publicação do aplicativo Android via \textit{Expo Application Services}. Para o TCC-III, está planejada a integração ponta a ponta da autenticação federada com \textit{Google}, a implementação das funcionalidades sociais (curtidas, indicação de presença e comentários) e a avaliação empírica de desempenho, custo e qualidade da extração estruturada.

\textbf{Palavras-chave}: Aplicativo Móvel. Computação \textit{Serverless}. Extração de Informação. Modelos de Linguagem. \textit{Web Crawling}.
\end{Resumo}

\begin{Abstract}
The discovery of cultural, sports, gastronomic, and community events in the northern coast of Santa Catarina faces a chronic fragmentation problem: each municipality, cultural foundation, and independent producer publishes its programming on separate websites, social networks, and printed materials, without structural standardization or regional aggregation. This undergraduate thesis presents the development of \textit{Rolê SC}, a mobile event aggregator application for the municipalities of Itajaí, Camboriú, Balneário Camboriú, Porto Belo, Itapema, and Bombinhas, supported by a serverless pipeline that automates the collection and structuring of this information. The ecosystem consists of four independent components deployed on \textit{Amazon Web Services}: a \textit{web crawler} based on \textit{Playwright} that collects pages from official websites; a natural language processor that uses the \textit{Llama 3.1 8B Instruct} model, served by the \textit{Groq} platform, to extract structured events or to discard pages that do not correspond to events; a \textit{REST API} developed in \textit{TypeScript} on the \textit{Fastify} framework, authenticated by \textit{JWT}; and an Android mobile application built with \textit{React Native} and \textit{Expo}. The work also covers the business model study, the analysis of direct and indirect competitors, requirements gathering, and system documentation. In this first stage (TCC-II), the collection and extraction pipeline was implemented for the five target municipalities, along with the flows for registration, email verification, authentication, password recovery and change, profile editing, and avatar upload, as well as the publication of the Android application through \textit{Expo Application Services}. For TCC-III, the end-to-end integration of federated authentication with \textit{Google}, the implementation of social features (likes, attendance indication, and comments), and the empirical evaluation of performance, cost, and structured extraction quality are planned.

\textbf{Keywords}: Mobile Application. \textit{Serverless} Computing. Information Extraction. Language Models. \textit{Web Crawling}.
\end{Abstract}


\pdfbookmark[0]{\listfigurename}{lof}
\listoffigures*
\cleardoublepage

\pdfbookmark[0]{\listtablename}{lot}
\listoftables*
\cleardoublepage


%\pdfbookmark[0]{\listofequacaoname}{loe}
%\listofequacao*
%\cleardoublepage

% --- Lista de siglas e abreviaturas ---
\begin{siglas}
  \item[ABNT] Associação Brasileira de Normas Técnicas
  \item[API] \textit{Application Programming Interface}
  \item[APK] \textit{Android Package}
  \item[AWS] \textit{Amazon Web Services}
  \item[CDN] \textit{Content Delivery Network}
  \item[CORS] \textit{Cross-Origin Resource Sharing}
  \item[CRUD] \textit{Create, Read, Update, Delete}
  \item[DOM] \textit{Document Object Model}
  \item[EAS] \textit{Expo Application Services}
  \item[HTML] \textit{HyperText Markup Language}
  \item[HTTP] \textit{HyperText Transfer Protocol}
  \item[HTTPS] \textit{HyperText Transfer Protocol Secure}
  \item[IAM] \textit{Identity and Access Management}
  \item[IaC] \textit{Infrastructure as Code}
  \item[JSON] \textit{JavaScript Object Notation}
  \item[JWT] \textit{JSON Web Token}
  \item[LGPD] Lei Geral de Proteção de Dados
  \item[LLM] \textit{Large Language Model}
  \item[MVP] \textit{Minimum Viable Product}
  \item[NLP] \textit{Natural Language Processing}
  \item[NoSQL] \textit{Not only SQL}
  \item[OAuth] \textit{Open Authorization}
  \item[OIDC] \textit{OpenID Connect}
  \item[OTA] \textit{Over-the-Air}
  \item[RF] Requisito Funcional
  \item[RN] Regra de Negócio
  \item[RNF] Requisito Não Funcional
  \item[REST] \textit{Representational State Transfer}
  \item[S3] \textit{Simple Storage Service}
  \item[SAM] \textit{Serverless Application Model}
  \item[SC] Santa Catarina
  \item[SDK] \textit{Software Development Kit}
  \item[SES] \textit{Simple Email Service}
  \item[SLA] \textit{Service Level Agreement}
  \item[SWOT] \textit{Strengths, Weaknesses, Opportunities, Threats}
  \item[TCC] Trabalho de Conclusão de Curso
  \item[TLS] \textit{Transport Layer Security}
  \item[TTL] \textit{Time to Live}
  \item[UC] Caso de Uso
  \item[UML] \textit{Unified Modeling Language}
  \item[URL] \textit{Uniform Resource Locator}
  \item[UX] \textit{User Experience}
\end{siglas}


% ---
\pdfbookmark[0]{\contentsname}{toc}
\tableofcontents*
\cleardoublepage
% ---

